\section{Additional Third-Party Batteries}
\label{sec:appendix-thirdparty}

This appendix reports the learner-family, attack-budget, and second-environment
batteries summarized in Section~\ref{sec:third-party}.

\subsection{Learner Boundary}

The same contrast attack ($\delta{=}2$) was evaluated on the third-party PPO
pipeline with action masking and sudden retirement.  It raised post-removal
violations in one of three paired runs, from 1178 to 2439
(McNemar exact $p=2.6\times10^{-149}$).  The first-violation episode changed
from 6 to 0.  The other two runs were unchanged ($p=1.00$ and $0.96$).  In a
separate ten-run shield-on-only isolation
battery, the run indexed by configured seed 101 was positive
($0.219\rightarrow0.754$ at retirement; McNemar exact
$p=3.7\times10^{-113}$), while the other nine entered a high-violation ceiling
regime (clean $0.728$--$0.760$) and were either protective or unchanged.
Figure~\ref{fig:ppo-isolation} reports the full battery.  A ten-run full-phase
PPO battery likewise placed every clean policy near the retirement ceiling
($0.73$--$0.77$) and yielded no positive outcome on the primary retirement
metric.  Its final metric captured a different effect: all ten poisoned
policies remained unsafe ($0.49$--$0.77$), whereas nine clean policies recovered
to approximately $0.22$ under continued learning.

A ten-run shield-on-only battery on off-policy soft actor-critic (SAC)
provided a complementary boundary.  Three runs were positive
($0.207\rightarrow0.561$, $0.218\rightarrow0.513$, and
$0.479\rightarrow0.954$; paired McNemar exact $p<10^{-40}$ for each), three
were protective, and four were unchanged.  All ten clean SAC policies had low
retirement disagreement rather than occupying the ceiling regime.  Four
extension runs generated on a second host were non-positive.  The fixed
attack therefore affected individual PPO and SAC training realizations without
producing reliable success within either family.

\begin{figure*}[t]
\centering
\includegraphics[width=0.85\textwidth]{fig6_ppo_isolation.pdf}
\caption{PPO retirement-boundary isolation on the third-party lifecycle
(shield-on-only contrast poisoning, $\delta{=}2$, paired runs indexed by
run indices 101--110,
sudden retirement).  At-retirement violation fraction of the raw policy for
clean and poisoned runs.  The run indexed by 101 is positive
($0.219\rightarrow0.754$; McNemar exact $p=3.7\times10^{-113}$); the other
nine runs land in the unsafe ceiling regime (clean at-retirement
$0.728$--$0.760$) with no paired headroom, six protective by the paired test and
three unchanged.  The dashed line marks the locked decision boundary
(poisoned at least $0.15$ above clean with paired McNemar $p<0.01$).}
\label{fig:ppo-isolation}
\end{figure*}

\subsection{Attack Budget and State Dependence}

On the full-phase REINFORCE protocol, the effect varied with attack shape,
budget, and training realization.  In the realization indexed by configured
seed~2, $\delta{=}2$ raised post-removal violations from 1127 to 3762
($3.3\times$; McNemar exact $p$ below numerical reporting precision).
Increasing the budget to $\delta{=}10$ produced 3766 violations.  In the
realization indexed by seed~3, $\delta{=}2$ was unchanged (2482 vs.\ 2471;
$p=0.84$), whereas $\delta{=}10$ raised violations to 3700
($p=8.4\times10^{-138}$).  The realization indexed by seed~1 began near the
unsafe ceiling (75.7\%) and showed no added damage at either budget.  Shape
controls were also state-dependent.  The bias variant was protective in the
ceiling realization.  The risk variant was protective in two realizations and
raised violations from 1127 to 2457 in the realization indexed by seed~2.

An independent Linux-host batch with identical dependency versions contained
one low-clean-violation realization and two ceiling realizations.  The former
showed the contrast effect at retirement ($0.232\rightarrow0.501$) and after
removal (1157 $\rightarrow$ 2451; $p=4.6\times10^{-159}$).  The latter two were
unchanged at $\delta{=}2$ and protective at $\delta{=}10$.  The run index of
the low-clean-violation regime differed across hosts.  Repeating the clean
configuration also entered a different trajectory mode ($0.776$ rather than
$0.232$ at retirement).  Dose comparisons are therefore confined to the
locked local batch.  The independent batch shows that attackability follows
the realized retirement state more closely than the nominal seed identifier.
Figure~\ref{fig:dose-response} summarizes both panels.

\begin{figure*}[t]
\centering
\includegraphics[width=\textwidth]{fig7_dose_response.pdf}
\caption{Budget- and state-dependent effects on the full-phase REINFORCE
protocol (paired within one locked code version).  (a) Contrast
attack dose response; (b) attack-shape negative controls at $\delta{=}2$.
The run-index-2 realization saturates at $\delta{=}2$; the run-index-3
realization requires $\delta{=}10$; the ceiling realization shows no additional damage at
either budget, and bias/risk controls are protective or strictly weaker.}
\label{fig:dose-response}
\end{figure*}

\subsection{Environment Boundary}

On a second gridworld of the same third-party suite (surveillance, $N{=}6$,
$R{=}3$), a clean reproduction again satisfied the authors' lifecycle
ordering: shield retained
$0/0$, smooth $0.129$, sudden $0.246$, and no-shield $0.252$ (retained $<$
smooth $\ll$ sudden $\le$ no-shield).  The fixed contrast attack at
$\delta{=}2$ raised post-removal violations in one of three independent paired
runs ($0.000 \rightarrow 0.272$, McNemar exact $p=2.6\times10^{-82}$) and was
protective in the other two ($0.267 \rightarrow 0.000$,
$p=8.4\times10^{-81}$; $0.135 \rightarrow 0.038$,
$p=1.3\times10^{-15}$).  The lifecycle ordering extended to the second
environment, whereas the exploit effect remained environment- and
realization-dependent.

\subsection{Opportunity-Region Cross-Tabulation}

Figure~\ref{fig:two-by-two} plots the coarse two-by-two
cross-tabulation summarized in Section~\ref{sec:third-party}: across 53
locked paired runs spanning both scopes, all three learner families, and
both domains, 12 of 20 runs with clean at-retirement violation below $0.5$
were positive, whereas none of 33 ceiling runs was positive.

\begin{figure*}[t]
\centering
\includegraphics[width=0.72\textwidth]{fig9_mechanism_divergence.pdf}
\caption{Clean at-retirement raw-policy/shield disagreement marks an
opportunity region across the obstacle corpus (both scopes, all three learner
families).  Each point is one paired clean/poisoned run; filled markers meet
the positive-outcome rule, and hollow markers do not.
Filled markers all have clean disagreement $\le0.047$; the dashed lines mark
the fixed low/high disagreement bands ($<0.06$ and $\ge0.09$).  No
high-disagreement run is positive, while seven of 18 low-disagreement runs are
non-positive.}
\label{fig:mechanism}
\end{figure*}

\begin{figure*}[t]
\centering
\includegraphics[width=\textwidth]{fig8_2x2_susceptibility.pdf}
\caption{Two-by-two opportunity-region cross-tabulation across all locked
paired runs (both scopes, all three learner families, both domains).  Each
point is one paired clean/poisoned run; filled markers meet the positive rule
(poisoned at-retirement at least $0.15$ above clean at retirement with paired
McNemar exact $p<0.01$), hollow markers do not.  The dashed rule line is
$y=x+0.15$ and the vertical dashed line marks the clean at-retirement $0.5$
boundary: positive outcomes concentrate below $0.5$, and no ceiling run is
positive.}
\label{fig:two-by-two}
\end{figure*}

\subsection{Family-Stratified Counts and Band Chronology}

Family-stratified counts clarify the AUC reported in
Section~\ref{sec:third-party}.  REINFORCE had seven positive low-band runs and
13 non-positive high-band runs.  PPO had one positive low-band run and 19
non-positive high-band runs.  All ten SAC runs were low-band, with three
positive and seven non-positive outcomes.

The protocol chronology bounds the statistical interpretation.  The
$0.06/0.09$ bands were fixed after analysis of 46 existing obstacle rows and
before four SAC extension pairs with run indices 407--410.  All four
extensions were low-band and non-positive.  They strengthened the
non-sufficiency result without prospectively testing high-disagreement
exclusion.  The marker is also obstacle-scoped: two protective runs in the
second environment had clean disagreement $0.476$ and $0.157$ despite clean
violation below $0.5$.

